\documentclass[10pt]{article}
\usepackage[T1]{fontenc}
\usepackage[utf8]{inputenc}
\usepackage[french]{babel}
\usepackage{amsmath,amssymb,mathrsfs,mathpazo}
\usepackage{geometry}
\geometry{a4paper,margin=2cm}
\usepackage{graphicx,array,tabularx,tabularray,colortbl,multirow,multicol,booktabs,enumitem,pifont,ulem,fancybox,diagbox,cancel}
\usepackage{pstricks,pstricks-add,pst-plot,pst-tree,pst-node,pst-text,pst-func,pst-math,pst-bspline,pst-3dplot,pst-eucl,pst-solides3d}
\let\pscalculate\undefined
\usepackage{pst-bezier}
\usepackage{tikz,pgfplots,tkz-tab}
\usetikzlibrary{arrows,arrows.meta,patterns,calc,positioning,trees,decorations.markings,shapes.geometric,shapes.misc,shapes.arrows,intersections,decorations.pathmorphing,decorations.pathreplacing,backgrounds,shadows,babel}
\pgfplotsset{compat=1.18}
\usepgfplotslibrary{fillbetween}
\usepackage[np]{numprint}
\usepackage{listings,makecell,pas-tableur,twoopt}
\let\barre\undefined
\usepackage[pstricks]{tablvar}
\definecolor{darkgreen}{RGB}{0,100,0}
\providecommand{\red}{\color{red}}
\providecommand{\blue}{\color{blue}}
\providecommand{\green}{\color{green}}
\providecommand{\black}{\color{black}}
\providecommand{\white}{\color{white}}
\providecommand{\gray}{\color{gray}}
\providecommand{\cyan}{\color{cyan}}
\providecommand{\magenta}{\color{magenta}}
\providecommand{\yellow}{\color{yellow}}
\providecommand{\violet}{\color{violet}}
\providecommand{\orange}{\color{orange}}
\providecommand{\euro}{\texteuro}
\providecommand{\EUR}{\texteuro}
\providecommand{\ieme}{\textsuperscript{e}}
\providecommand{\degres}{\ensuremath{^\circ}}
\providecommand{\degree}{\ensuremath{^\circ}}
\providecommand{\vv}[1]{\overrightarrow{#1}}
\providecommand{\Coord}[1]{#1}
\providecommand{\decosix}{}
\providecommand{\starredbullet}{\ensuremath{\bullet}}
\pagestyle{empty}
\setlength{\parindent}{0pt}
\usepackage{adjustbox}\DeclareUnicodeCharacter{2236}{:}\DeclareUnicodeCharacter{2190}{\ensuremath{\leftarrow}}
\providecommand{\R}{\mathbb{R}}
\providecommand{\N}{\mathbb{N}}
\providecommand{\Z}{\mathbb{Z}}
\providecommand{\Q}{\mathbb{Q}}
\providecommand{\D}{\mathbb{D}}
\providecommand{\C}{\mathbb{C}}
\providecommand{\e}{{\mathrm{e}}}
\providecommand{\vect}[1]{\overrightarrow{#1}}
\providecommand{\vectt}[1]{\overrightarrow{\mathrm{#1}}}
\providecommand{\vvt}[1]{\overrightarrow{\mathrm{#1}}}
\providecommand{\np}[2][]{{#2~#1}}
\providecommand{\vv}[1]{{\overrightarrow{#1}}}
\renewcommand{\d}{{\mathrm{d}}}
\providecommand{\degres}{{^\circ}}
\providecommand{\degre}{{^\circ}}
\providecommand{\blue}{\color{blue}}
\providecommand{\red}{\color{red}}
\providecommand{\enskip}{\quad}
\providecommand{\llbracket}{\lbrack\!\lbrack}
\providecommand{\rrbracket}{\rbrack\!\rbrack}
\providecommand{\euro}{{\text{€}}}
\providecommand{\no}{{\text{n°}}}
\providecommand{\interval}[2]{{[#1;#2]}}
\providecommand{\intervFO}[2]{{[#1;#2[}}
\providecommand{\dotfill}{\ldots}

\providecommand{\ds}{\displaystyle}
\providecommand{\textding}[1]{#1}
\providecommand{\Oij}{\(\mathrm{O};\vec{\imath},\vec{\jmath}\)}
\providecommand{\Oijk}{\(\mathrm{O};\vec{\imath},\vec{\jmath},\vec{k}\)}
\providecommand{\Ouv}{\(\mathrm{O};\vec{u},\vec{v}\)}
\providecommand{\intervFF}[2]{\ensuremath{[#1;#2]}}
\renewcommand{\euro}{\ifmmode\text{\texteuro}\else\texteuro\fi}
\DeclareUnicodeCharacter{00B0}{\ifmmode^{\circ}\else\textdegree\fi}
\DeclareUnicodeCharacter{20AC}{\ifmmode\text{\texteuro}\else\texteuro\fi}
% Énoncés : Éducation nationale. Transcription et corrigés : APMEP, Denis Vergès et contributeurs. Découpage et classement : CoopMaths, Sébastien Lozano, Éric Elter et contributeurs. Adaptation de la présentation : Maths Explorer.
% CC BY-SA : https://coopmaths.fr/www/about/
% https://coopmaths.fr/alea/static/e3c/2020/tex/e3c_2020_00_sujet22_1.tex
\begin{document}
\noindent\textit{APMEP / CoopMaths. CC BY-SA. Présentation adaptée par Maths Explorer.}\par\medskip
\medskip

\emph{Ce QCM comprend cinq questions indépendantes.}

\emph{Pour chacune des questions, une seule des quatre réponses proposées est correcte.}

\emph{Pour chaque question, indiquer le numéro de la question et recopier sur la copie la lettre correspondant à la réponse choisie.}

\emph{ \textbf{Aucune justification n’est demandée mais il peut être nécessaire d’effectuer des recherches} au brouillon pour aider à déterminer votre réponse.}

\emph{Chaque réponse correcte rapporte $1$ point. \textbf{Une réponse incorrecte ou une question sans réponse n’apporte ni ne retire de point}.}

\medskip

\textbf{Question 1}

\medskip

Sur la figure ci-dessous, nous avons tracé dans un repère orthonormé la courbe représentative $\mathcal{C}$ d’une fonction $f$ dérivable sur $\R$ et la tangente à $\mathcal{C}$ au point d’abscisse 4. Cette tangente est représentée par la droite $\mathcal{D}$. On note $f'$ la fonction dérivée de la fonction $f$ .

\begin{center}
\psset{labelFontSize=\scriptstyle,showorigin=false}
\begin{pspicture}(-3.5,-3.8)(6,4.5)
\multido{\n=-2+1}{8}{\psline[linewidth=0.75pt,linecolor=lightgray](\n,-3.5)(\n,4.2)}
\multido{\n=-3+1}{8}{\psline[linewidth=0.75pt,linecolor=lightgray](-2.5,\n)(5.3,\n)}
\psaxes[linewidth=0.95pt,]{->}(0,0)(-2.9,-3)(5.5,4.4)
\psdot[dotstyle=+,dotscale =1.4,dotangle=45](4,-1)
\uput[ur](4,-1){A} 
\psplot[linewidth=1.15pt,linecolor=blue,plotpoints=2000]{-1}{5}{x 2 sub 2 exp 0.5 neg  mul 1 add}
\psplot[linewidth=1.15pt,linecolor=red,plotpoints=2000]{1.32}{5}{x 2 neg   mul 7 add}
\uput[dl](0,0){O}\uput[d](0.5,0.7){\blue $\mathcal{C}$}\uput[d](2.25,3.5){\red $\mathcal{D}$}
\end{pspicture}
\end{center}

Le réel $f'(4)$ est égal à :

\medskip

\begin{tabularx}{\linewidth}{*{4}{X}}
\textbf{a.~~} $-1$ &\textbf{b.~~} $-2$&\textbf{c.~~} $7$& \textbf{d.~~} $1$.
\end{tabularx}

\medskip

\textbf{Question 2}

\medskip

Soit $f$ la fonction définie sur $\R$ par $ f(x) = x^3 - 2x^2 + 1$. On admet que $f$ est une fonction dérivable sur $\R$. Dans un repère, une équation de la tangente à la courbe représentative de la fonction $f$ au point d’abscisse 1 est :

\medskip

\begin{tabularx}{\linewidth}{*{4}{X}}
\textbf{a.~~} $y=-1$ &\textbf{b.~~} $y=-x$&\textbf{c.~~} $y=-x+1$& \textbf{d.~~} $y=x$.
\end{tabularx}

\medskip
\textbf{Question 3}
\medskip

Pour tout réel $x$ , $\dfrac{\e^x\times\e^{-3x}}{\e^{-x}}$ est égal à :

\medskip

\begin{tabularx}{\linewidth}{*{4}{X}}
\textbf{a.~~} $ \e^{-x}$ &\textbf{b.~~} $\e^{3x}$&\textbf{c.~~}$\e^{-3x}$& \textbf{d.~~} $\e^{x}$.
\end{tabularx}

\medskip

\textbf{Question 4}

\medskip

Soit $f$ une fonction polynôme du second degré dont la courbe représentative dans un repère orthonormé est donnée ci-dessous.

\begin{center}
\psset{labelFontSize=\scriptstyle,showorigin=false}
\begin{pspicture}(-4,-5)(3,4)
\multido{\n=-3+1}{6}{\psline[linewidth=0.75pt,linecolor=lightgray](\n,-4.5)(\n,3.2)}
\multido{\n=-4+1}{8}{\psline[linewidth=0.75pt,linecolor=lightgray](-3.5,\n)(2.5,\n)}
\psaxes[linewidth=0.95pt,]{->}(0,0)(-3.5,-4.8)(2.5,3.4)


\psplot[linewidth=1.15pt,linecolor=blue,plotpoints=5000]{-2.48}{1.5}{x  2 exp 2 mul x 2  mul add 4 sub}

\uput[dl](0,0){O}
\end{pspicture}
\end{center}

Pour tout réel $x$, une expression de $f(x)$ est :

\medskip

\begin{tabularx}{\linewidth}{*{4}{X}}
\textbf{a.~~}$ f(x)=x^2+x-2$&\textbf{b.~~}$f(x)=-x^2-4$&\textbf{c.~~}$f(x)=2x^2+2x-4$&\textbf{d.~~}$f(x)=-3x^2-3x+6$
\end{tabularx}

\medskip

\textbf{Question 5}

\medskip

L’ensemble  $\mathcal{S}$ des solutions de l’inéquation d’inconnue $x\in\R$ : $-x^2-2x+8>0$ est :

\medskip

\begin{tabularx}{\linewidth}{*{4}{X}}
\textbf{a.~~} $\mathcal{S} = [-4~;~2]$&\textbf{b.~~} $\mathcal{S}=]-4~;~2[$&\textbf{c.~~}\footnotesize $\mathcal{S} = ]-\infty~;~-4]\cup]2~;~+\infty[$&\textbf{d.~~}  $\mathcal{S} = \{-4~;~2\}$\
\end{tabularx}

\vspace{0,5cm}

\end{document}
